\begin{ouverture}[probabilites]{11}{Probabilités, évènements et simulations}{Partie 3 · Données et probabilités}
{Issues · Évènements · Réunion · Intersection · Tableaux · Fréquences}
\objectif{Décrire des évènements avec le vocabulaire des ensembles.}
\objectif{Utiliser la relation $P(A\cup B)+P(A\cap B)=P(A)+P(B)$.}
\objectif{Dénombrer les issues de deux épreuves indépendantes et interpréter une simulation.}
\prerequis{Fractions, probabilité d’un évènement et de son contraire.}
\end{ouverture}
\pagestyle{otMain}\essentiel\fontsize{11}{15}\selectfont
\essSection{1}{Évènements et opérations}
\begin{definition}Une issue est un résultat possible d’une expérience aléatoire. Un évènement est un ensemble d’issues. L’évènement contraire de $A$, noté $\overline A$, contient les issues qui ne réalisent pas $A$. L’évènement impossible correspond à l’ensemble vide $\varnothing$.\end{definition}
\begin{definition}L’intersection $A\cap B$ (« $A$ et $B$ ») contient les issues communes à $A$ et $B$. La réunion $A\cup B$ (« $A$ ou $B$ ») contient les issues de $A$, de $B$, ou des deux. $A$ et $B$ sont incompatibles si $A\cap B=\varnothing$.\end{definition}
\begin{propriete}Si les issues sont équiprobables, $P(A)$ est le nombre d’issues favorables divisé par le nombre total d’issues. De plus $P(\overline A)=1-P(A)$ et
\[P(A\cup B)+P(A\cap B)=P(A)+P(B).\]\end{propriete}
\begin{exemple}Avec un dé équilibré, $A$ : « pair » $=\{2;4;6\}$ et $B$ : « au moins 4 » $=\{4;5;6\}$. $A\cap B=\{4;6\}$ et $A\cup B=\{2;4;5;6\}$ : on vérifie $4/6+2/6=3/6+3/6$.\end{exemple}
\begin{attention}Le « ou » des mathématiques est inclusif. « $A$ ou $B$, mais pas les deux » se dit « ou exclusif » ; sa probabilité vaut $P(A\cup B)-P(A\cap B)$.\end{attention}
\clearpage
\essSection{2}{Deux épreuves et simulations}
\begin{savoirfaire}[Dénombrer]Pour deux épreuves indépendantes, un tableau à double entrée ou un arbre de dénombrement liste toutes les issues. Si chaque épreuve a des issues équiprobables, les couples obtenus le sont aussi : deux dés donnent $6\times6=36$ couples équiprobables.\end{savoirfaire}
\begin{exemple}Deux pièces équilibrées lancées indépendamment donnent quatre issues équiprobables : PP, PF, FP, FF. L’évènement « exactement une face F » contient PF et FP, donc a pour probabilité $2/4=1/2$.\end{exemple}
\begin{savoirfaire}[Numéroter pour compter]Dans un sac de 3 boules rouges et 2 bleues, numéroter les boules $R_1,R_2,R_3,B_1,B_2$. Avec remise, deux tirages donnent $5\times5=25$ couples équiprobables, dont $3\times3=9$ formés de deux rouges.\end{savoirfaire}
\begin{propriete}Simuler une expérience, c’est la reproduire par un dispositif ou un programme qui respecte le modèle, par exemple un nombre entier aléatoire de 1 à 6 pour un dé. Quand le nombre de répétitions augmente, la fréquence observée se stabilise généralement autour de la probabilité. Cela ne prédit ni le prochain résultat ni une égalité exacte.\end{propriete}
\clearpage\rubrique{Découverte}
\begin{blocLivre}{ex-11-1}\begin{exercice}[1]{Un dé}
On lance un dé équilibré à six faces numérotées de 1 à 6. Quelle est la probabilité d’obtenir un nombre pair ?
\end{exercice}\end{blocLivre}
\begin{blocLivre}{ex-11-2}\begin{exercice}[1]{Des boules colorées}
Un sac contient 5 boules rouges, 3 bleues et 2 vertes, indiscernables au toucher. On tire une boule au hasard. Donner les probabilités de rouge et de non-rouge.
\end{exercice}\end{blocLivre}
\begin{blocLivre}{ex-11-3}\begin{exercice}[1]{Écrire des évènements}
On lance un dé équilibré. $A$ est l’évènement « obtenir un nombre pair » et $B$ l’évènement « obtenir un multiple de 3 ». Écrire $A$, $B$, $A\cap B$, $A\cup B$ et $\overline A$ sous forme d’ensembles.
\end{exercice}\end{blocLivre}
\begin{blocLivre}{ex-11-4}\begin{exercice}[1]{Impossible ou certain}
Sur un dé classique, donner les probabilités d’obtenir 9 et d’obtenir moins de 7. Comment note-t-on l’évènement impossible ?
\end{exercice}\end{blocLivre}
\begin{blocLivre}{ex-11-5}\begin{exercice}[1]{Une fréquence expérimentale}
On obtient 46 fois face au cours de 100 lancers d’une pièce. Donner la fréquence de face. Est-ce une preuve que la pièce est déséquilibrée ?
\end{exercice}\end{blocLivre}

\clearpage\rubrique{Entraînement}
\begin{blocLivre}{ex-11-6}\begin{exercice}[2]{La relation de la réunion}
Avec les évènements de l’exercice 3, calculer $P(A)$, $P(B)$ et $P(A\cap B)$, puis en déduire $P(A\cup B)$ par la relation du cours. Contrôler par un dénombrement.
\end{exercice}\end{blocLivre}
\begin{blocLivre}{ex-11-7}\begin{exercice}[2]{Deux pièces}
Deux pièces équilibrées sont lancées indépendamment. Compléter l’arbre de dénombrement et calculer la probabilité d’obtenir exactement une face F, puis au moins une face F.
\label{fig-ex-11-7}\figureGeo{proba-pieces}
\end{exercice}\end{blocLivre}
\begin{blocLivre}{ex-11-8}\begin{exercice}[2]{La somme de deux dés}
Deux dés équilibrés sont lancés indépendamment. Utiliser le tableau des sommes pour calculer la probabilité d’une somme égale à 7.
\label{fig-ex-11-8}\figureGeo{proba-des}
\end{exercice}\end{blocLivre}
\begin{blocLivre}{ex-11-9}\begin{exercice}[2]{Un « ou » exclusif}
On lance un dé équilibré. $A$ : « pair » et $B$ : « au moins 4 ». Calculer la probabilité de l’évènement « $A$ ou $B$, mais pas les deux ».
\end{exercice}\end{blocLivre}
\begin{blocLivre}{ex-11-10}\begin{exercice}[2]{Au moins un six}
Deux dés équilibrés sont lancés. On note $A$ : « le premier dé donne 6 » et $B$ : « le second dé donne 6 ». Calculer la probabilité d’obtenir au moins un six.
\label{fig-ex-11-10}\figureGeo{proba-des}
\end{exercice}\end{blocLivre}

\clearpage\rubrique{Problèmes}
\begin{blocLivre}{ex-11-11}\begin{exercice}[3]{Deux tirages avec remise}
Un sac contient 3 boules rouges $R_1,R_2,R_3$ et 2 bleues $B_1,B_2$. On tire une boule, on la remet, puis on en tire une seconde. Utiliser le tableau pour calculer la probabilité de deux rouges, puis de deux boules de même couleur.
\label{fig-ex-11-11}\figureGeo{proba-tableau-boules}
\end{exercice}\end{blocLivre}
\begin{blocLivre}{ex-11-12}\begin{exercice}[3]{Un jeu à points}
On paie 2 points pour lancer un dé équilibré. Le 6 rapporte 8 points, le 5 rapporte 4 points, les autres 0. Sur un grand nombre de parties, le gain brut moyen théorique est la somme des gains multipliés par leurs probabilités. Le jeu équilibre-t-il mise et gain moyen ?
\end{exercice}\end{blocLivre}
\begin{blocLivre}{ex-11-13}\begin{exercice}[3]{Un effectif inconnu}
Un sac contient 6 rouges et $b$ bleues seulement. La probabilité de tirer rouge au hasard est $2/5$. Trouver $b$.
\end{exercice}\end{blocLivre}
\begin{blocLivre}{ex-11-14}\begin{exercice}[3]{Un code à deux chiffres}
Un générateur choisit indépendamment et au hasard deux chiffres parmi 0, 1, 2, 3, avec répétition possible. Quelle est la probabilité que ses deux chiffres soient distincts ? Que leur somme soit paire ou que le premier chiffre soit 0 ?
\end{exercice}\end{blocLivre}
\begin{blocLivre}{ex-11-15}\begin{exercice}[3]{Estimer sans promettre}
Un programme simule un évènement de probabilité $1/3$ à l’aide de nombres aléatoires. Il obtient 327 succès en 1000 essais, puis 3354 en 10000 essais. Comparer les écarts entre fréquences et probabilité. Peut-on garantir que tout essai supplémentaire rapproche la fréquence de $1/3$ ?
\end{exercice}\end{blocLivre}

\clearpage
